\documentclass[11pt]{article}
\usepackage[margin=1in]{geometry}
\usepackage{amsmath,amssymb,amsthm,mathtools}
\usepackage{booktabs,array}
\usepackage[T1]{fontenc}
\usepackage{lmodern,microtype}
\usepackage{xcolor,hyperref}
\hypersetup{colorlinks=true,linkcolor=blue!40!black,citecolor=blue!40!black,
  urlcolor=blue!50!black,
  pdftitle={The adaptive plain-moment estimate: supplementary proof},
  pdfauthor={Zhongpai Gao}}
\newtheorem{theorem}{Theorem}[section]
\newtheorem{proposition}[theorem]{Proposition}
\theoremstyle{remark}
\newtheorem{remark}[theorem]{Remark}
\newcommand{\Q}{\mathbb Q}
\newcommand{\OO}{\mathcal O_F}
\newcommand{\N}{\mathrm N}
\newcommand{\Rea}{\operatorname{Re}}
\newcommand{\one}{\mathbf1}
\newcommand{\norm}[1]{\left\lVert#1\right\rVert}
\title{The adaptive plain-moment estimate:\\supplementary proof}
\author{Zhongpai Gao\\\small gaozhongpai@gmail.com}
\date{9 October 2026}
\begin{document}
\maketitle
\begin{abstract}
We extend the parameter range of the plain fourth-moment estimate in
OpenAI's September 2026 preprint from $\kappa\ge3/4$ to
$\kappa\ge13/18$, assuming the corresponding global zero bound for
finite-order Hecke characters. We verify the comparison inequalities,
the admission of every recursive child, and the order of the mesh,
height and induction choices. We then give the resulting witness count.
The unchanged arithmetic identities and analytic estimates are inputs
from the cited, unrefereed preprint. This supplement does not assert
an independent proof of those inputs or perform the subsequent
zero-free-region continuation.
\end{abstract}

\section{The estimate and its hypotheses}\label{adp:statement}

Put $F=\Q(\sqrt{-3})$. We use the primary generators, sextic residue
symbols and fixed finite set of excluded primes $\mathcal S$ of
\cite[Sections 4 and 14]{OAI26}. For an integral ideal or its primary
generator $n$, write $q_n=\N n$; for an element $k$, write
$q_k=|k|^2$. Sums over $n$ or $l$ use the primary representatives of
ideals prime to $\mathcal S$. The character $\chi_n(k)$ is the
zero-extended sextic residue symbol. In particular, a displayed sixth
power is one on units and is zero on nonunits.

Let $\beta_*$ be the supremum of $1/2$ and of the real parts of zeros
in $1/2\le\Rea s\le1$ of primitive finite-order Hecke $L$-functions
over $F$. Principal poles are excluded. This is the global family
quantity in \cite[Section 2, (2.1)]{OAI26}.

Fix a finite group $\Theta$ of finite-order ray characters supported
on $\mathcal S$, containing the fixed twists, the characters used to
separate reciprocity phases, and $n\mapsto\chi_n(-1)$. At scale
$Z\ge2$, the rows and their characters are
\[
 0<q_k\ll Z^m,\qquad \psi_k(n)=\tau(n)\chi_n(k),\qquad M=m+q .
\]
Here $\tau$ is fixed within the row sum. The union of the displayed
moving prime supports in $\tau$ has norm at most $Z^q$. It is measured
\emph{before} cancellation or reduction of exponents modulo six.
Every factor retains its original zeros. An additional squarefree
mask $\mathfrak R$, of norm at most $Z^{B_0}$ for fixed $B_0$, may
depend on frozen parameters but is fixed within a row sum and is
common to both plain sums and every prime factor.

For smooth profiles supported on fixed annuli, put
\begin{align}
 S_\psi(n;W)&=Z^{-n/2}\sum_l\psi(l)\one_{(l,\mathfrak R)=1}
                                      W(q_l/Z^n),\label{adp:plain-def}\\
 Q_{\psi,i}&=Z^{-z_i/2}\sum_{p\ {\rm prime}}
       \psi(p)\one_{(p,\mathfrak R)=1}\nu_i(p)W_i(q_p/Z^{z_i}),
       \qquad
 Q_\psi=\prod_iQ_{\psi,i},\quad z=\sum_i z_i.\label{adp:slot-def}
\end{align}
Each $\nu_i$ is a fixed finite linear combination of members of
$\Theta$, independent of the row and all other prime choices.
The underlying prime windows are disjoint before imposing any zeros
or masks. A factor of length zero has bounded scale and is estimated
by absolute counting.
Let $\mathcal R_0$ contain the rows with nonprincipal primitive
inducing character. If $z>0$, let $\mathcal R_z$ contain the rows
whose primitive inducing character is not induced by a member of
$\Theta$. Membership refers to the inducing character, without
removing the presentation's zeros.

\begin{theorem}[Extended plain moment]\label{adp:plain}
Retain the arithmetic and profile hypotheses of
\cite[Lemma 18.1]{OAI26}. Suppose
\[
 \frac{13}{18}\le\kappa\le\frac34,\qquad
 \beta_*\le s_\kappa:=\frac{1+\kappa}{2}.
\]
For every $\varepsilon>0$ and prescribed bounded ranges of nonnegative
lengths, there is a mesh $\eta>0$, chosen before the fixed number of
prime factors, such that, whenever $z_i\le\eta$,
\begin{equation}\label{adp:moment}
 \sum_{k\in\mathcal R_z}
 \left|S_{\psi_k}(n_1;W_1)S_{\psi_k}(n_2;W_2)Q_{\psi_k}\right|^2
 \ll Z^{M+\varepsilon}.
\end{equation}
For $z=0$, the lengths $n_1,n_2$ are unrestricted in their prescribed
bounded ranges, and no hypothesis on $\beta_*$ is required.
For $z>0$, assume
\begin{equation}\label{adp:admission}
 n_1+n_2+6\kappa z\le M .
\end{equation}
The estimate retains the finite smooth-seminorm and polynomial
height bounds of the original lemma. Their orders and constants
may depend on the fixed data and number of prime factors, but are
uniform in rows, moving labels, masks and the numerical heights.
The mesh is uniform for $\kappa\in[13/18,3/4]$.
\end{theorem}

The global hypothesis is essential: the two Poisson transforms create
new finite-order characters, which need not have the same rightmost
zero as the original row. We prove the extension by checking every
use of the lower bound on $\kappa$ in \cite[Section 18]{OAI26}.
Its exact reflection, finite correlations, common-coefficient identity
and masked lattice cancellation are imported with all their original
hypotheses. The argument below changes their numerical admissions;
it does not replace their coefficients by arbitrary bounded weights.
